\section{Related Work}
\label{sec:related-work}

In the following, we summarize related work regarding causal knowledge 
graphs, approaches for causal question answering, and approaches that apply reinforcement learning to reasoning tasks on knowledge graphs.

\paragraph{Causal Knowledge Graphs}
ConceptNet~\cite{Speer2017ConceptNet} is a general knowledge graph (KG) consisting of 36 relations between natural language terms, including a \textit{Causes} relation. CauseNet~\cite{Heindorf2020Causenet} and Cause Effect Graph~\cite{Li2020CauseEffectGraph} specifically focus on causal relations extracted via linguistic patterns from web sources like Wikipedia and ClueWeb12. ATOMIC~\cite{Sap2019ATOMIC} focuses on inferential knowledge of commonsense reasoning in everyday life. It consists of \textit{``If-Event-Then-X''} relations based on social interactions or real-world events. ATOMIC$^{20}_{20}$~\cite{Hwang2021COMET} selects relations from ATOMIC and ConceptNet to create an improved graph while adding more relations via crowdsourcing. Instead, \citet{West2022Symbolic} automate the curation of causal relations by clever prompting of a language model. Finally, CSKG~\cite{Ilievski2021CSKG} builds a consolidated graph combining seven KGs, including ConceptNet and ATOMIC. We use CauseNet because most of the other KGs are smaller and less focused on causal relations, e.g., CauseNet contains many more causal relationships than ConceptNet~\cite{Heindorf2020Causenet}. Future work may involve combining multiple KGs.

\paragraph{Causal Question Answering}
As of now, only few approaches tackle the causal question answering task. Most of them focus on binary questions, i.e., questions such as ``\textit{Does X cause Y?}'' which expect a ``yes'' or ``no'' answer. 
\citet{KayeshCausalTransfer2020} model the task as a transfer learning approach. They extract cause-effect pairs from news articles via causal cue words. Subsequently, they transform the pairs into sentences of the form ``\textit{X may cause Y}'' and use them to finetune BERT~\cite{DevlinBert2019}. Similarly, \citet{HassanzadeshCausalQA2019} employ large-scale text mining to answer binary causal questions introducing multiple unsupervised approaches ranging from string matching to embeddings computed via BERT. \citet{SharpCausalQAEmbeddings2016} consider multiple-choice questions of the form ``\textit{What causes X?}''. First, they mine cause-effect pairs from Wikipedia via syntactic patterns and train an embedding model to capture the semantics between them. At inference time, they compute the embedding similarity between the question and each answer candidate. \citet{DalalCausalQAEnhancing2021, DalalCausalQAISWC2021}
combine a language model with CauseNet~\cite{Heindorf2020Causenet}.
Given a question, they apply string matching to extract relevant causal relations from CauseNet. Subsequently, they provide the question with the causal relations as additional context to a language model. As other language model-based approaches, too, they cannot produce verifiable answers. None of them selects relevant paths in a causality graph via reinforcement learning.

\paragraph{Knowledge Graph Reasoning with Reinforcement Learning}
In recent years, reinforcement learning on knowledge graphs has been successfully applied to link prediction~\cite{Das2018Minerva}, fact-checking~\cite{Xiong2017DeePpath}, or question answering~\cite{Qiu2020Stepwise}. Given a source and a target entity, DeepPath~\cite{Xiong2017DeePpath} learns to find paths between them. The training of DeepPath involves two steps. First, it is trained via supervised learning and afterward via REINFORCE~\cite{Williams1992REINFORCE} policy gradients. During inference time, the paths are used to predict links between entities or check the validity of triples. Subsequently, MINERVA~\cite{Das2018Minerva} improves on DeepPath by introducing an LSTM~\cite{Hochreiter1997LSTM} into the policy network to account for the path history. Moreover, MINERVA does not require knowledge of the target entity and is trained end-to-end without supervision at the start. \citet{Lin2020RewardShaping} propose two improvements for MINERVA. First, they apply reward shaping by scoring the paths with a pre-trained KG embedding model~\cite{Dettmers2018ConvE} to reduce the problem of sparse rewards. Second, they introduce a technique called action dropout, which randomly disables edges at each step. Action dropout serves as additional regularization and helps the agent learn diverse paths. M-Walk~\cite{Shen2018MWalk} applies model-based reinforcement learning techniques. Like AlphaZero~\cite{Silver2018AlphaZero}, M-Walk applies Monte Carlo Tree Search (MCTS) as a policy improvement operator. Thus, at each step, M-Walk applies MCTS to produce trajectories of an improved policy and subsequently trains the current policy to imitate the improved one. GaussianPath~\cite{Wan2021GaussianPath} takes a Bayesian view of the problem and represents each entity by a Gaussian distribution to better model uncertainty.

Previous approaches for reinforcement learning on KGs~\cite{Xiong2017DeePpath, Das2018Minerva, Qiu2020Stepwise} were designed for KGs with multiple relation types and used structural KG embeddings. In contrast, we tailor our approach to causality graphs where nodes are noun phrases, edges are all of the same type (“causes”), and each edge has provenance data associated with it. This results in the following key differences:
(1)~
Different action space:
As edge types do not provide any learning signal, our action space encompasses all entities in the causality graph.
(2)~
Different entity encoding: 
We use text embeddings instead of structural KG embeddings because causality graphs are text-centric, e.g., nodes are noun phrases.
(3)~
Different action encoding:
An action corresponds to the traversal of an edge that leads to an adjacent node. We encode an action as $\textbf{a}=[\textbf{s};\textbf{e}]$ where $\textbf{s}$ is the sentence embedding of he sentence that the edge was extracted from and $\textbf{e}$ is the embedding of the adjacent entity.
(4)~
Different RL algorithm:
We train our agent with the more advanced Synchronous Advantage Actor-Critic (A2C) algorithm instead of the REINFORCE algorithm used in previous works.
